\documentclass[11pt,a4paper]{article}
\usepackage[margin=22mm]{geometry}
\usepackage{amsmath,amssymb}
\usepackage{fontspec}
\usepackage{polyglossia}
\setdefaultlanguage{russian}
\setmainfont{DejaVu Serif}
\usepackage[unicode,hidelinks]{hyperref}

\allowdisplaybreaks[1]
\raggedright
\setlength{\parskip}{0.35em}
\newcommand{\dd}{\mathop{}\!\mathrm{d}}
\newcommand{\avg}[1]{\langle #1\rangle}
\newcommand{\ket}[1]{\lvert #1\rangle}
\newcommand{\bra}[1]{\langle #1\rvert}
\newcommand{\braket}[2]{\langle #1\mid #2\rangle}
\DeclareMathOperator{\Var}{Var}

\newcommand{\solutiontitle}[4]{%
  \hypersetup{pdftitle={Квантовая механика 1. Решения ДЗ №#1}}%
  \begin{center}
    {\small Квантовая механика 1 \enspace\textperiodcentered\enspace Семестр 1\par}
    \vspace{1ex}
    {\LARGE\bfseries Решения домашнего задания №#1\par}
    \vspace{0.5ex}
    {\large #2\par}
  \end{center}
  \noindent{\small Условия приведены в файле «#3», #4.}\par\medskip
}

\begin{document}
\solutiontitle{11}{Гармонический осциллятор и лестничные операторы}{ДЗ2026(1)\_11.pdf}{с.~1}

Положим \(\ell=\sqrt{\hbar/(m\omega)}\), \(Q=x/\ell\),
\(P=\ell p/\hbar\). Тогда \([Q,P]=iI\) и
\[
 a=\frac{Q+iP}{\sqrt2},\qquad
 a^\dagger=\frac{Q-iP}{\sqrt2},\qquad
 [a,a^\dagger]=I,\qquad H=\hbar\omega(N+\tfrac12),\quad N=a^\dagger a .
\]
Действительно, \(2N=Q^2+P^2+i[Q,P]=Q^2+P^2-I\).
В фоковском базисе \(N\ket n=n\ket n\).
Коммутатор \([a,a^\dagger]=I\) показывает, что \(a,a^\dagger\)
меняют \(n\) на \(-1,+1\) соответственно.
Квадраты норм \(a\ket n\) и \(a^\dagger\ket n\) равны
\(n,n+1\), поэтому при стандартном выборе фаз
\(a\ket n=\sqrt n\ket{n-1}\) и
\(a^\dagger\ket n=\sqrt{n+1}\ket{n+1}\).

\section{Эволюция заданной суперпозиции}

При \(P=-i\partial_Q\) вакуум решает \((Q+\partial_Q)\varphi_0=0\).
Нормировка и два последовательных действия
\((Q-\partial_Q)/\sqrt2\) дают
\[
 \varphi_0=\pi^{-1/4}e^{-Q^2/2},\qquad
 \varphi_1=\sqrt2\,Q\varphi_0,\qquad
 \varphi_2=\frac{2Q^2-1}{\sqrt2}\varphi_0 .
\]
Из \(Q^2-Q=\tfrac12+(2Q^2-1)/2-Q\) следует
\[
 \Psi(Q,0)=A\pi^{1/4}
       \left(\tfrac12\varphi_0-\tfrac1{\sqrt2}\varphi_1
                                 +\tfrac1{\sqrt2}\varphi_2\right).
\]
Норма в квадрате равна \(|A|^2\sqrt\pi(1/4+1/2+1/2)\).
Берём \(A=2/(\sqrt5\,\pi^{1/4})\).
Коэффициенты \(c_0,c_1,c_2\) равны \((1,-\sqrt2,\sqrt2)/\sqrt5\).

\textbf{(a)} Умножая каждую компоненту на её стационарную фазу,
получаем
\[
 \Psi(Q,t)=\frac{1}{\sqrt5}\bigl(
 \varphi_0 e^{-i\omega t/2}-\sqrt2\,\varphi_1e^{-3i\omega t/2}
                    +\sqrt2\,\varphi_2e^{-5i\omega t/2}\bigr).
\]

\textbf{(b)} По вероятностям \(|c_n|^2\) имеем
\(\avg{N}=0/5+2/5+4/5=6/5\), откуда
\(\avg{E}=17\hbar\omega/10\).

\textbf{(c)} Две соседние пары дают среднее значение
\[
 \avg{a(t)}=e^{-i\omega t}(c_0^*c_1+\sqrt2\,c_1^*c_2)
           =-\frac{3\sqrt2}{5}e^{-i\omega t}.
\]
Так как \(Q=(a+a^\dagger)/\sqrt2\),
\(P=(a-a^\dagger)/(i\sqrt2)\), получаем
\[
 \avg{Q(t)}=-\frac65\cos\omega t,\qquad
 \avg{P(t)}=\frac65\sin\omega t .
\]

\textbf{(d)} Из \(Q^2=N+1/2+(a^2+(a^\dagger)^2)/2\) и
\(\avg{a^2}_{t=0}=\sqrt2\,c_0^*c_2=2/5\) следует
\(\avg{Q^2}_{t=0}=17/10+2/5=21/10\).
Поскольку \(m\omega^2 x^2/2=(\hbar\omega/2)Q^2\),
\(\avg{U}_{t=0}=21\hbar\omega/20\).

\section{Матрицы в фоковском базисе}

Нумеруем строки и столбцы с нуля. Действие на \(\ket n\)
даёт \(a_{mn}=\sqrt n\,\delta_{m,n-1}\) и
\((a^\dagger)_{mn}=\sqrt{n+1}\,\delta_{m,n+1}\).
Поэтому первые четыре строки и столбца имеют вид
\[
 a=\begin{pmatrix}
 0&1&0&0\\0&0&\sqrt2&0\\0&0&0&\sqrt3\\0&0&0&0
 \end{pmatrix},\qquad
 a^\dagger=\begin{pmatrix}
 0&0&0&0\\1&0&0&0\\0&\sqrt2&0&0\\0&0&\sqrt3&0
 \end{pmatrix}.
\]
Линейные комбинации дают соответствующие блоки
\[
 Q=\frac1{\sqrt2}\begin{pmatrix}
 0&1&0&0\\1&0&\sqrt2&0\\0&\sqrt2&0&\sqrt3\\0&0&\sqrt3&0
 \end{pmatrix},\qquad
 P=\frac1{i\sqrt2}\begin{pmatrix}
 0&1&0&0\\-1&0&\sqrt2&0\\0&-\sqrt2&0&\sqrt3\\0&0&-\sqrt3&0
 \end{pmatrix}.
\]
Для квадратов сначала раскрываем полные операторы,
\(Q^2=N+1/2+(a^2+(a^\dagger)^2)/2\) и
\(P^2=N+1/2-(a^2+(a^\dagger)^2)/2\).
Так как \(a^2\ket n=\sqrt{n(n-1)}\ket{n-2}\) для \(n\ge2\)
и \(a^2\ket n=0\) для \(n=0,1\), блоки равны
\[
 Q^2=\frac12\begin{pmatrix}
 1&0&\sqrt2&0\\0&3&0&\sqrt6\\\sqrt2&0&5&0\\0&\sqrt6&0&7
 \end{pmatrix},\qquad
 P^2=\frac12\begin{pmatrix}
 1&0&-\sqrt2&0\\0&3&0&-\sqrt6\\-\sqrt2&0&5&0\\0&-\sqrt6&0&7
 \end{pmatrix}.
\]
Формулы через \(a,a^\dagger,N\) задают также все элементы за
пределами выписанных блоков.

\section{Коммутаторы лестничных операторов}

Сохраняя порядок множителей в правиле Лейбница, имеем
\begin{align*}
 [a,N]&=[a,a^\dagger]a+a^\dagger[a,a]=a,\\
 [a^\dagger,N]&=[a^\dagger,a^\dagger]a+a^\dagger[a^\dagger,a]=-a^\dagger,\\
 [a^\dagger,a^2]&=[a^\dagger,a]a+a[a^\dagger,a]=-2a,\\
 [a^\dagger,N^2]&=[a^\dagger,N]N+N[a^\dagger,N]
                 =-a^\dagger N-Na^\dagger=-a^\dagger(2N+I).
\end{align*}
В последней строке использовано \(Na^\dagger=a^\dagger(N+I)\).

\section{Движение средних в произвольном состоянии}

Пусть \(\sum_n |C_n|^2=1\) и состояние имеет конечное
\(\avg{N}\), что достаточно для существования нужных средних.
При действии \(a\) на \(\ket{n+1}\) в скалярном
произведении остаётся компонента \(\bra n\).
Разность энергий соседних уровней всегда \(\hbar\omega\), откуда
\[
 z=\sum_{n\ge0}\sqrt{n+1}C_n^*C_{n+1},\qquad
 \avg{a(t)}=z e^{-i\omega t},\qquad
 \avg{a^\dagger(t)}=z^* e^{i\omega t}.
\]
Сумма для \(z\) сходится абсолютно по неравенству Коши --- Буняковского
с множителями \((|C_n|)\) и \((\sqrt{n+1}|C_{n+1}|)\).
Положим \(A_0=\sqrt2\operatorname{Re}z=\avg{Q(0)}\),
\(B_0=\sqrt2\operatorname{Im}z=\avg{P(0)}\). Тогда
\[
 \avg{Q(t)}=A_0\cos\omega t+B_0\sin\omega t,\qquad
 \avg{P(t)}=B_0\cos\omega t-A_0\sin\omega t .
\]
Таким образом константы \((A,B,C,D)\) из условия равны
\((A_0,B_0,B_0,-A_0)\).
Классические безразмерные переменные подчиняются
\(\dot Q=\omega P,\ \dot P=-\omega Q\).
Решение этой системы совпадает по форме с найденным движением
средних, где начальные данные равны \(\avg{Q(0)},\avg{P(0)}\).

\section{Лестничные операторы в представлении Гейзенберга}

Формулы задачи 3 дают \([H,a]=-\hbar\omega a\) и
\([H,a^\dagger]=\hbar\omega a^\dagger\).
Поэтому \(\dot a_H=-i\omega a_H,\ \dot a^\dagger_H=i\omega a^\dagger_H\)
с начальными операторами \(a,a^\dagger\). Решение
\[
 a_H(t)=e^{-i\omega t}a,\qquad
 a_H^\dagger(t)=e^{i\omega t}a^\dagger .
\]
Тот же фазовый множитель можно проверить действием
\(e^{iHt/\hbar}a e^{-iHt/\hbar}\) на каждый \(\ket n\).

\end{document}
